\subsubsection{Linear Operators}

\paragraph{$L^p$ Space}

\begin{definition}[ space]\textbf{$L^p$}\\
Let $(X,\Sigma, \mu)$ be a measure space and $1\leq p\leq \infty$. The space $L^p(X,\mu)$ is the set of all measurable functions from $X$ to $\mathbb{R}$ or $\mathbb{C}$ such that $\int_X |f(x)|^p d\mu < \infty$.

\begin{itemize}
\item Note $L^p$ is a \textit{normed space} with norm:

\begin{equation}
\|f\|_{L_p}=\left(\int_X |f|^p d\mu\right)^{1/p}
\end{equation}


$p$ can be the symbol $\infty$ with the norm defined as:

\begin{equation}
\|f\|_{L_{\infty}} = \inf\{C\in \mathbb{R}_{\geq 0} : |f(x)|<C, \forall x\}
\end{equation}


\item It is \textit{complete}, hence it is a Banach Space
\end{itemize}

\end{definition}\paragraph{Linear Operators}

\begin{definition}[Linear Operators]Let $T$ be a map betwen two normed space $X$ and $Y$. $T$ is linear if:
\begin{equation}
T(x+\lambda y) = T(x) +\lambda T(y)
\end{equation}
$\forall x,y \in X, \lambda \in K$

\end{definition}If we consider linear functions $f:\mathbb{R}\to \mathbb{R}$, it seems continuity is free since we can always write $T(x)=kx$ hence:
\begin{equation}
\|T(x)-T(x_0)\|=\|k(x-x_0)\|=k\|x-x_0\|
\end{equation}
so we can easily find a delta ball of radius $\frac{\epsilon}{k}$ given an epsilon ball such that the delta ball is mapped into the epsilon ball. More generally we can define bounded operator to make the above logic work:

\begin{definition}[Bounded Operators]Let $T$ be a map betwen two normed space $X$ and $Y$. $T$ is bounded if $\exists C>0$:
\begin{equation}
\|T(x)\|_Y\leq C\|x\|_X
\end{equation}
$\forall x \in X$

\end{definition}\begin{proposition}[Continuous and Bounded are equivalent for Linear Operator]The following are equivalent:

\begin{enumerate}
\item If $X$ and $Y$ are normed spaces and $T:X\to Y$ is a linear map which is continuous from $X$ to $Y$
\item $T$ is bounded
\end{enumerate}

\end{proposition}\begin{proof}\begin{itemize}
\item Suppose $T$ is \textbf{linear and bounded}, $\forall x_0 \in X$, $\forall \epsilon > 0$, consider the open ball around $x_0$ with radius $\epsilon$ $B_{\epsilon}^o$. Then $\forall x\in B_{\epsilon}^o$:
\begin{equation}
\|T(x)-T(x_0)\| &= \|T(x-x_0)\| \\
&\leq C\|x-x_0\|
\end{equation}
Now to make $C\|x -x_0\|<\epsilon$, we need $\|x -x_0\|<\frac{\epsilon}{C}$. Hence $\delta(\epsilon)=\frac{\epsilon}{C}$


\item Suppose $T$ is \textbf{linear and continous}. Then in particular $T$ is continous at $x=0$. We can find a $\delta >0$ such that $\|x\| < \delta \Rightarrow \|T(x)\| < 1$. Now to control arbitrary $x\in X$, we can rewrite $x$ as
\begin{equation}
x=\left(\frac{2}{\delta }\right)\|x\|\left(\frac{\delta x}{2\|x\|}\right)=C\|x\|x'
\end{equation}
Note $\|x'\|=\frac{\delta }{2}\|x\|<\delta$ hence $\|T(x')\|\leq 1$. By linearity $T(x)=T(C\|x\|x')=C\|x\|T(x')$. Thus:
\begin{equation}
\|T(x)\|_Y=C\|x\|\|T(x')\|\leq C\|x\|
\end{equation}
\end{itemize}

\end{proof}It turns out all finite-dimensional linear operator are bounded. However, infinite-dimensional linear operator can easily be unbounded.